\subsection{域的特征指数。完全域}

设 K 为一个域。K 的特征指数指的是这样一个整数：若 K 的特征为 0，它等于 1；若特征非零，它等于 K 的特征。

命题 4. --- 设 K 是特征指数为 q 的域。对每个整数 \(f \geqslant 0\)，映射 \(x \mapsto x^{q^{f}}\) 是 K 到它的一个子域上的同构（记作 \(K^{q^{f}}\)）。

当 \(q \neq 1\) 时这由定理 2 得出，而当 \(q = 1\) 时是平凡的。

同样地，可以把命题 2 推广到 A 是特征指数为 q 的域的情形，q = 1 的情形是平凡的。

定义 4. --- 一个特征指数为 q 的域 K 称为完全域，如果 \(K^{q} = K\) 。当 \(K^{q} \neq K\) 时，称 K 为非完全域。

根据这一定义，一个域是完全的，如果它的特征为 0，或者它是定义 2 意义下特征为 \(p \neq 0\) 的完全环。如果 \(\mathbf{K}\) 是一个特征为 \(p \neq 0\) 的域，且 \((\hat{\mathbf{K}}, u)\) 是 \(\mathbf{K}\) 的一个完全闭包，则根据命题 3（V，第 6 页），\(\hat{\mathbf{K}}\) 是一个域，且 \(u\) 是 \(\mathbf{K}\) 到 \(\hat{\mathbf{K}}\) 的某个子域上的一个同构。通常把 \(\mathbf{K}\) 与它在 \(u\) 下的像在 \(\hat{\mathbf{K}}\) 中等同起来，于是我们有 \(\hat{\mathbf{K}} = \mathbf{K}^{p^{-\infty}}\)（命题 3）。

设 K 为特征为 0 的域；则 K 的特征指数为 1。按约定，记号 \(x^{q^{-f}}\) 与 \(S^{q^{-f}}\) 分别表示 x 和 S（其中 x 是 K 的元素，S 是 K 的子集）。特别地，我们令 \(K^{q^{-\infty}} = K\)，并且约定取 K 的完全闭包为 K 本身。

命题 5. --- 若 K 是特征为 0 的域，或者是有限域*或代数闭域*，则它是完全域。特别地，每个素域都是完全域。

假设 K 的特征为 \(p \neq 0\) 。若 K 为有限域，则子域 \(K^{p}\) 与 K 具有相同的基数，从而 \(K^{p} = K\)。* 若 K 是代数闭的，则多项式 \(X^{p} - a\) 对每个 \(a \in K\)（V，第 20 页，定义 1）在 K 中有一个根 x，因此 \(x^{p} = a\)，于是 \(K^{p} = K\)。* 最后，素域的特征为 0 或有限。

设 \(K_{0}\) 是特征为 \(p \neq 0\) 的域，\(\mathbf{K} = \mathbf{K}_{0}(\mathbf{X})\) 是不定元 X 上 \(K_{0}\) 上的有理分式域。则 K 是非完全域，因为不存在 K 的元素 \(u(\mathbf{X})/v(\mathbf{X})\)（u，v 是 \(K_{0}[\mathbf{X}]\) 中的多项式），使得 \((u(\mathbf{X})/v(\mathbf{X}))^{p} = \mathbf{X}\)。把该关系写成 \(u(\mathbf{X})^{p} = \mathbf{X}v(\mathbf{X})^{p}\) 的形式并比较两边的次数即可看出这一点。

